\scanpage{60}\mspage{59}
Donc, en \deleted{\unsure{de|te}}termes d'une structure portée par un ens ordonné $M$ (cf remarque page 58, on trouve ceci. On a sur \deleted{\uline{Définition}} $M$ une relation $R$ (``compatibilité) qui est réflexive symétrique (donc équivaut à la donnée de $\mathcal S\subset\mathcal P_2(M)$), \deleted{\uline{impliquant}}\inserted{et telle que} $R(x,y)$ implique $R(x',y')$ pour $x'\leq x$, $y'\leq y$. \deleted{\unsure{comp stable} par $\leq$.} Cela permet, à peu de choses près, de reconstituer $(\mathcal F,\leq)$, en \unsure{tant qu'}ens de parties de $M$. Pour l'expliciter,

Quelques terminologies

``\uline{Composante}'' \uline{irréductible} \inserted{(ou strate maximale)} d'une figure $F$ : \deleted{\unsure{o|a}}strate maximale de $F$ ($\exists$ tjrs si l'ens $\overline F$ des strates de $F$ est fini \inserted{(on dit alors que $F$ est une \uline{figure finie})} et $F\ne\varnothing_{\mathcal F}$).

Figure \deleted{de type}\inserted{de type fini} fini : qui est borne sup. de \unsure{l'ens|la réu} de ses ``composantes'' irréductibles (\unsure{exemple|exemples} : toute figure finie est de t.f.), ou encore, toute figure \inserted{$F$} qui est \deleted{de la forme $\operatorname{Sup}X$}\inserted{borne supr. d'un \unsure{nb} \uline{fini} de} figures irréductibles. \deleted{Ceci} (NB Dans les cas les plus importants, toute figure est finie, et a fortiori de type fini.)

Ceci dit, si $\mathcal F_{\mathrm{t.f.}}$ désigne l'ens figures de type fini, on récupère $\overline{\mathcal F_{\mathrm{t.f}}}$, l'ens des parties \unsure{correspondantes|correspondants} de $M$, par la seule connaissance de $(M,\leq,R)$,
